Large language model agents increasingly act on the world through external tools,
and the Model Context Protocol (MCP) has become the dominant standard for
connecting them. In under two years MCP has gone from a single vendor's proposal
to an ecosystem with public registries listing tens of thousands of servers; the
official registry alone contains \Nframe{} unique servers at the time of our study.
When an agent selects a tool and an MCP server mishandles the request, by returning
a malformed result, accepting ill-typed arguments, or failing to start, the agent
does not merely lose a capability, and may proceed on wrong output with full
confidence. The reliability of this server layer is therefore a reliability
property of every agent built on top of it.

Yet we know remarkably little about how these servers actually behave, because most
published empirical studies of the MCP ecosystem examine servers \emph{without
running them}, and those that do run them look for vulnerabilities. Existing work scans source code and metadata for health and
maintainability, crawls registries to characterize scale and supply-chain
structure, or mines GitHub issues to taxonomize faults reported by users. These
static lenses answer valuable questions, but they cannot answer a prior and more
basic one: \emph{if you install these servers and speak the protocol to
them, do they work, and do they follow the specification?} The protocol maintainers
ship a conformance test suite, but it is built for SDK maintainers to run in
continuous integration, and to our knowledge it has not been used to measure the
deployed ecosystem. The SDKs are tested; the seventeen thousand servers built on them
are not.

This paper addresses that gap with an execution-based conformance study of the
MCP server ecosystem. Starting from the official registry, we execute a complete
census of every self-contained, locally-installable server, launch each in an
isolated sandbox, and drive it through a conformance harness that exercises the
initialization lifecycle, tool listing, declared-schema validity, and behavior
under three inputs that arise in ordinary agent sessions: a call to an unknown
tool, a call with a wrong-typed argument, and a malformed protocol frame. We grade
each server against the protocol version it negotiates, and we take care to
separate genuine specification requirements (RFC-2119 \textsf{MUST}/\textsf{SHOULD})
from behavior the specification merely illustrates by example.

Only \HandshakeRateCorr{} of installable servers start and complete a protocol
handshake. The rest fail in ways we classify into a startup-failure taxonomy, from
install errors to silent exits to hangs. Among servers that do respond,
\ErrAsResultRate{} report a call to a non-existent tool as a ``successful''
tool-execution error (\texttt{isError: true}) rather than the JSON-RPC protocol error
the specification illustrates. We trace this near-universal divergence to the default
behavior of the official SDKs, which is inherited by the servers built on
them. The result bears on a live debate, in which an open proposal would
reclassify unknown-tool failures to match what the SDKs already do, arguing the
change aligns the specification with practice~\cite{sep2145}. Our census says how far
practice has already gone, which decides whether such a change ratifies the ecosystem
or redirects it. Separately, \SilentExecRate{} of responding servers execute a tool on
an argument their own declared schema rejects and return a well-formed result. In one
case a code-scanning tool asked to scan an integer replied \texttt{CLEAN}, with
neither an error code nor an \texttt{isError} flag a client could act on. Unlike the
error-reporting question, this behaviour violates an explicit requirement that servers
validate all tool inputs. A further \MalformedDiesRate{} fail to recover from a malformed frame and
answer nothing further, although JSON-RPC 2.0 requires a reply to every request.

We organise the study around five research questions:

\begin{itemize}
  \item \textbf{RQ1 (Runnability).} What fraction of registry-published servers can a
    client install and start, and when they do not start, why not?
  \item \textbf{RQ2 (Lifecycle conformance).} Of the servers that respond, how many
    correctly implement the initialization lifecycle, tool listing, and declared-schema
    validity?
  \item \textbf{RQ3 (Input enforcement).} How many servers execute a tool on an
    argument their own declared schema rejects, and what does the client receive when
    they do?
  \item \textbf{RQ4 (Specification versus practice).} Where the specification and
    deployed behaviour disagree on error reporting, how large is the divergence, what
    causes it, and which side is moving?
  \item \textbf{RQ5 (Correlates).} Do runnability and conformance vary with the
    package registry a server ships on?
\end{itemize}

This paper makes four contributions:
\begin{itemize}
  \item \textbf{A protocol-level conformance measurement of the MCP ecosystem by
    execution} (C1): where prior ecosystem studies read code and metadata, we install
    and speak the protocol to a complete census of the eligible registry, testing
    behaviour against inputs rather than inspecting declarations, and releasing an
    open harness and dataset so the measurement is fully reproducible.
  \item \textbf{A two-dimensional taxonomy} (C2) of startup-failure classes and
    conformance-violation classes, each tied to a concrete agent-facing consequence.
  \item \textbf{A root-cause account of the spec--practice divergence} (C3): by
    attributing each server to its SDK, we show the dominant divergence is
    institutionalized in the official SDKs rather than independently reinvented.
  \item \textbf{Open infrastructure} (C4): a conformance harness that extends the
    official suite's scenarios to arbitrary deployed servers, plus the released
    measurement dataset, enabling the ecosystem to monitor its own conformance.
\end{itemize}
